\documentclass[letterpaper,11pt]{article}
\usepackage[margin=1in]{geometry}
\usepackage{times}
\usepackage[hyphens]{url}
\usepackage{graphicx}
\usepackage{amsmath,amssymb}
\usepackage{booktabs}
\usepackage{multirow}
\usepackage{xcolor}
\usepackage{hyperref}
\usepackage{natbib}
\usepackage{enumitem}

\setcounter{secnumdepth}{2}

\title{Reproducibility Budgets for LLM Evaluation:\\A Generalizability Theory Analysis of Variance Sources\\Across Benchmarks, Models, and Evaluation Metrics}

\author{
Anonymous Authors
}

\begin{document}
\maketitle

% ============================================================
% ABSTRACT
% ============================================================
\begin{abstract}
When LLM evaluation results vary across reproductions, practitioners cannot determine whether the culprit is prompt phrasing, random seed, numerical precision, or item selection---let alone how these factors interact.
We conduct the first multi-benchmark Generalizability Theory (G-theory) study of LLM evaluation, crossing infrastructure (numerical precision) with methodology (prompt template, random seed, temperature, option ordering) and measurement (test items) factors across five benchmarks (MMLU, ARC-Challenge, HellaSwag, GSM8K, MATH) and seven 7--9B parameter models.
Three findings emerge.
\textbf{(1)}~Item selection is the largest single variance source on knowledge-retrieval benchmarks (42.1\% on MMLU) but \emph{model$\times$item interaction} dominates on reasoning benchmarks (53.7\% on MATH), yielding a monotonically decreasing gradient of item dominance from knowledge retrieval to multi-step reasoning (permutation test $p=0.008$).
\textbf{(2)}~Numerical precision$\times$methodology interactions are negligible under binary scoring ($<$0.8\%), confirming aggregate separability on a single GPU architecture.
\textbf{(3)}~G-theory requires 3.6--14.4$\times$ more items than IRT-based approaches to reach $G\geq0.80$, because it accounts for prompt, seed, and ordering dimensions invisible to IRT; at IRT-recommended item counts, G-theory reliability ranges from 0.312 (MATH) to 0.529 (MMLU).
Decision studies (D-studies) yield benchmark-specific reproducibility budgets: 340 items suffice for $G\geq0.80$ on MMLU, but 780 are needed for MATH.
We release an open-source G-theory pipeline and D-study budget calculator.
\end{abstract}

% ============================================================
% 1. INTRODUCTION
% ============================================================
\section{Introduction}

Researchers debating whether to standardize prompts, random seeds, or numerical precision for LLM evaluation implicitly assume that their findings will generalize across benchmarks.
A growing body of work has documented substantial variance from individual evaluation factors: prompt formatting can swing accuracy by up to 76\% \citep{sclar2024quantifying}, numerical precision introduces up to 9\% standard deviation \citep{yuan2025llminference}, and different random seeds produce 80 out of 1,000 unique completions even at temperature zero \citep{thinkingmachines2025}.
Yet each study isolates at most one or two sources, leaving practitioners without a unified account of how these factors interact---or whether their relative importance depends on the benchmark under study.

Generalizability Theory \citep[G-theory;][]{shavelson1991generalizability,brennan2001generalizability} provides a principled framework for simultaneously decomposing multiple variance sources and projecting reliability under alternative measurement designs.
\citet{mesting2026reliability} adapted G-theory to LLM evaluation, decomposing five methodological facets on a single benchmark and finding judge choice (43.8\%) as the dominant source.
However, that study treated infrastructure as fixed and analyzed only one benchmark, leaving open whether variance structure is benchmark-conditional.
Meanwhile, IRT-based approaches \citep{lalor2019learning,polo2024tinybenchmarks} have addressed efficient item subsampling but from a single-facet perspective that does not account for prompt, seed, or ordering variance.

We conduct a multi-benchmark G-theory study crossing numerical precision with four methodological facets (prompt template, random seed, temperature, option ordering) across five benchmarks and seven models.
By treating model as a crossed quasi-random facet, we reveal that item difficulty rankings are substantially model-specific: the item$_{\text{id}}$ component drops from 61\% in single-model analyses to 42\% when model is crossed, while model$\times$item interaction emerges at 34\%.
Crucially, this variance structure is \emph{not} stable across benchmarks---item dominance monotonically decreases from knowledge-retrieval benchmarks (MMLU: 42.1\%) to multi-step reasoning benchmarks (MATH: 14.2\%), while model$\times$item interaction increases correspondingly (33.6\%$\to$53.7\%).
This cross-benchmark gradient, confirmed by permutation test ($p=0.008$), renders any ``universal reproducibility recipe'' illusory.

Our contributions are:
\begin{enumerate}
    \item \textbf{Cross-benchmark variance gradient.} We provide the first empirical demonstration that the dominant variance source in LLM evaluation shifts systematically from item difficulty (knowledge retrieval) to model$\times$item interaction (multi-step reasoning), with a 27.9 percentage-point spread confirmed by permutation-tested monotonicity ($p=0.008$). D-study predictions validated out-of-sample ($\Delta=0.004$).
    \item \textbf{Precision separability (metric-qualified).} Numerical precision$\times$methodology interactions are negligible under binary scoring ($<$0.8\%), providing empirical support for the implicit separability assumption in isolated-factor studies. Under text-exact-match scoring, precision$\times$temperature reaches 4.15\%, consistent with floating-point non-associativity.
    \item \textbf{IRT vs.\ G-theory: the cost of multi-facet reliability.} G-theory requires 3.6--14.4$\times$ more items than IRT for $G\geq0.80$, because it accounts for prompt$\times$item, seed$\times$item, and ordering$\times$item dimensions that IRT ignores. This quantifies the hidden cost of evaluation dimensions invisible to classical item analysis.
    \item \textbf{Practitioner toolkit.} We release an open-source G-theory variance decomposition pipeline, D-study budget calculator, and benchmark-specific reproducibility budget lookup tables.
\end{enumerate}

% ============================================================
% 2. RELATED WORK
% ============================================================
\section{Related Work}

\paragraph{Variance in ML evaluation.}
\citet{bouthillier2021accounting} established factorial variance decomposition for ML benchmarks via ANOVA, crossing hyperparameters, data splits, and random seeds.
\citet{henderson2018deep} documented reproducibility challenges in deep reinforcement learning, finding that implementation details often dominate algorithmic differences.
\citet{dodge2020finetuning} and \citet{sellam2022multberts} quantified fine-tuning instability from random seeds and data ordering in BERT-scale models.
\citet{dehghani2021benchmark} argued that benchmark selection itself introduces a ``lottery'' effect on model rankings.

\paragraph{LLM evaluation variance.}
\citet{sclar2024quantifying} demonstrated that prompt formatting alone can swing accuracy by up to 76\%, though \citet{hua2025emnlp} showed that much apparent prompt sensitivity is an evaluation artifact of log-likelihood scoring.
\citet{yuan2025llminference} documented up to 9\% accuracy variation from numerical precision without factorial attribution.
\citet{brittlebench2026} decomposed prompt sensitivity at approximately 50\% variance on fixed hardware.
\citet{messina2026background} introduced a reference-calibration framework for quantifying background randomness at temperature zero.
\citet{kunievsky2025decomposition} proposed an alternative decomposition into intent, articulation, and model uncertainty.

\paragraph{G-theory in LLM evaluation.}
\citet{mesting2026reliability} adapted Generalizability Theory to LLM evaluation, decomposing five methodological facets (item, prompt, judge, temperature, seed) and finding judge choice as the dominant variance source (43.8\%).
Our work extends this by: (a) crossing infrastructure (numerical precision) with methodology in a unified framework, (b) treating model as a crossed facet rather than fixed, (c) analyzing five benchmarks to reveal the cross-benchmark variance gradient, and (d) fixing the judge to isolate infrastructure$\times$methodology interactions---\citeauthor{mesting2026reliability}'s 43.8\% judge variance would otherwise dominate.

\paragraph{IRT approaches to LLM evaluation.}
Item Response Theory has been applied to NLP evaluation by \citet{lalor2019learning}, \citet{rodriguez2021evaluation}, and \citet{vania2021comparing}.
\citet{polo2024tinybenchmarks} used IRT-based subsampling to recommend $\sim$100 items for efficient benchmarking.
These approaches optimize item selection along a single facet (item difficulty/discrimination), without accounting for prompt, seed, or ordering variance.
Our D-study budgets are directly comparable to IRT recommendations, quantifying the additional cost of multi-facet reliability.

% ============================================================
% 3. METHOD
% ============================================================
\section{Method}

\subsection{Generalizability Theory Background}

G-theory \citep{shavelson1991generalizability,brennan2001generalizability} partitions observed score variance into components attributable to each facet and their interactions in a crossed measurement design.
For an observation $X_{pifso}$ on person $p$ (here, model), item $i$, prompt template $f$, seed $s$, and ordering $o$, the linear model is:
\begin{align}
X_{pifso} &= \mu + \alpha_p + \beta_i + \gamma_f + \delta_s + \epsilon_o \notag\\
&\quad + (\alpha\beta)_{pi} + (\alpha\gamma)_{pf} + \cdots + e_{pifso}
\label{eq:gtheory}
\end{align}
where each Greek letter denotes a random effect for the corresponding facet, parenthesized terms are interactions, and $e$ is residual.
Variance components $\hat{\sigma}^2$ are estimated via Henderson's Method I (expected mean squares), which is unbiased for balanced designs \citep{brennan2001generalizability}.

A \textbf{generalizability coefficient} summarizes the ratio of universe-score variance to observed-score variance:
\begin{equation}
G = \frac{\hat{\sigma}^2_p}{\hat{\sigma}^2_p + \sum_j \hat{\sigma}^2_{\delta_j}/n_j}
\label{eq:gcoeff}
\end{equation}
where the sum runs over all error components weighted by design sizes $n_j$.
A \textbf{Decision study} (D-study) projects $G$ under hypothetical designs (e.g., varying numbers of items, prompts, seeds), enabling practitioners to determine the measurement budget required for a target reliability.

\paragraph{Binary data.}
Our primary metric (binary correct/incorrect) violates G-theory's normality assumption.
We address this through a four-part defense: (a) matched-design simulations showing binary and continuous Henderson I recover equivalent dominant-source rankings ($\rho=0.732$ vs.\ $0.741$ at $n=200$); (b) complementary GLMM analysis on a subsample; (c) Bernoulli-adjusted decomposition confirming percentage decomposition is inherently normalized by $p(1-p)$; and (d) logprob continuous-metric analysis showing gradient preservation (Section~\ref{sec:robustness}).
Null simulations confirm that observed variance structure is not a Bernoulli artifact: simulated item$_{\text{id}}$ under the null is 0.08\% versus observed 42.1\%.

\subsection{Experimental Design}

\paragraph{Benchmarks.}
We study five benchmarks spanning knowledge retrieval to multi-step reasoning: MMLU \citep{hendrycks2021measuring}, ARC-Challenge \citep{clark2018think}, HellaSwag \citep{zellers2019hellaswag}, GSM8K \citep{cobbe2021training}, and MATH \citep{hendrycks2021math}.
Each benchmark contributes 200 randomly sampled items.
This range was chosen to test whether variance structure is benchmark-conditional, as hypothesized from psychometric theory: knowledge-retrieval items should show more stable difficulty rankings across models than reasoning items requiring model-specific solution strategies.

\paragraph{Models.}
Seven open-weight models in the 7--9B parameter range: Qwen3-8B, Llama-3.1-8B, Gemma-2-9B, Mistral-7B-v0.3, Yi-1.5-9B, OLMo-2-7B, and InternLM2.5-7B (3 of 5 benchmarks).
Models were purposively sampled from major architecture families to maximize architectural diversity within a fixed parameter range.
All inference uses vLLM 0.9.2 with generation-based evaluation following \citet{hua2025emnlp} to avoid log-likelihood scoring artifacts.

\paragraph{Facets.}
Each benchmark crosses the following facets:
\begin{itemize}
    \item \textbf{Numerical precision}: BF16, FP16 (2 levels for cross-benchmark; 3 levels including FP32 for MMLU-specific analyses).
    \item \textbf{Prompt template}: 6 semantically equivalent templates per benchmark.
    \item \textbf{Random seed}: 6 seeds controlling CUDA-level nondeterminism under default vLLM settings.
    \item \textbf{Option ordering}: 4 systematic rotations for multiple-choice benchmarks; 4 rephrasing variants for free-form benchmarks.
    \item \textbf{Temperature}: $T \in \{0.0, 0.7\}$ for cross-benchmark analyses.
\end{itemize}
Model is treated as a crossed quasi-random facet.
The total design yields approximately 2,016,000 item-level binary records across all phases.

\paragraph{Estimation.}
Variance components are estimated via Henderson's Method I with REML cross-validation.
REML-EM (500 iterations, convergence $\Delta < 0.003$ pp) confirms all Henderson I estimates within 0.47 pp (rank correlation $\rho = 1.0$), consistent with the near-balanced design.
Bootstrap 95\% confidence intervals use 1,000 item-level resamples; jackknife comparison confirms CI width ratio 0.96--1.04.

% ============================================================
% 4. RESULTS
% ============================================================
\section{Results}

\subsection{Cross-Benchmark Variance Decomposition}
\label{sec:crossbenchmark}

Table~\ref{tab:crossbenchmark} presents the Henderson I variance decomposition across all five benchmarks with model treated as a crossed facet.
The dominant variance source shifts systematically across benchmarks: item$_{\text{id}}$ accounts for 42.1\% of variance on MMLU but only 14.2\% on MATH, while model$\times$item interaction moves in the opposite direction, from 33.6\% (MMLU) to 53.7\% (MATH).

\begin{table*}[t]
\centering
\caption{Cross-benchmark variance decomposition (Henderson I, 7 models $\times$ BF16). Percentages denote proportion of total variance. $G_{\text{item}}$ is the generalizability coefficient for item scores across models. D-study reports minimum items for $G \geq 0.80$.}
\label{tab:crossbenchmark}
\small
\begin{tabular}{lccccc}
\toprule
\textbf{Component} & \textbf{MMLU} & \textbf{ARC} & \textbf{HellaSwag} & \textbf{GSM8K} & \textbf{MATH} \\
\midrule
item$_{\text{id}}$ (\%) & 42.1 & 35.8 & 27.4 & 19.6 & 14.2 \\
model$\times$item (\%) & 33.6 & 37.2 & 41.8 & 48.3 & 53.7 \\
model (\%) & 0.5 & 0.8 & 1.2 & 2.4 & 3.1 \\
prompt$\times$item (\%) & 2.8 & 2.6 & 2.9 & 1.4 & 0.8 \\
seed$\times$item (\%) & 2.5 & 2.1 & 1.8 & 1.2 & 0.6 \\
ordering$\times$item (\%) & 0.6 & 0.4 & 0.3 & 0.1 & 0.1 \\
residual (\%) & 19.5 & 21.3 & 24.6 & 26.8 & 27.4 \\
\midrule
$G_{\text{item}}$ & 0.820 & 0.776 & 0.714 & 0.653 & 0.612 \\
D-study items ($G \geq 0.80$) & 340 & 420 & 510 & 660 & 780 \\
\bottomrule
\end{tabular}
\end{table*}

This pattern admits a coherent interpretation grounded in task structure.
Knowledge-retrieval benchmarks (MMLU, ARC) test factual recall where item difficulty is relatively model-invariant---most models either ``know'' a fact or do not---yielding high item$_{\text{id}}$ and comparatively low model$\times$item interaction.
As tasks shift toward multi-step reasoning (HellaSwag$\to$GSM8K$\to$MATH), correctness increasingly depends on model-specific reasoning strategies: different models deploy different solution paths (e.g., algebraic manipulation vs.\ numerical substitution), and their success-failure patterns diverge at the item level.
This drives model$\times$item interaction upward while compressing item$_{\text{id}}$.

The model main effect is small across all benchmarks (0.5--3.1\%), indicating that aggregate accuracy differences between 7--9B models contribute minimally to total variance---item-level heterogeneity dominates.
Methodological facets (prompt, seed, ordering) collectively account for $<$6\% on all benchmarks, interacting primarily with items rather than contributing as main effects.
Residual variance increases from 19.5\% (MMLU) to 27.4\% (MATH), reflecting additional stochasticity in generation-based evaluation of reasoning tasks (chain-of-thought path dependence, format variability in free-form answers).

\subsection{Gradient Monotonicity}
\label{sec:gradient}

The cross-benchmark gradient in item$_{\text{id}}$ is not merely a trend---it is monotonically decreasing from knowledge retrieval to multi-step reasoning.
Table~\ref{tab:monotonicity} reports formal tests.

\begin{table}[t]
\centering
\caption{Gradient monotonicity tests. The permutation test assesses whether the observed ordering of item$_{\text{id}}$\% across five benchmarks could arise by chance (5,000 permutations).}
\label{tab:monotonicity}
\small
\begin{tabular}{llcl}
\toprule
\textbf{Test} & \textbf{Statistic} & \textbf{$p$} & \textbf{Conclusion} \\
\midrule
Permutation & $\tau = -1.0$ & 0.008 & Reject $H_0$ \\
 (item$_{\text{id}}$\% gradient) & (perfect monotone) & & at $\alpha = 0.01$ \\
\addlinespace
Within-MC & $\Delta = -6.3$ pp & 0.023 & Significant \\
(MMLU$\to$ARC) & & & within format \\
\addlinespace
Within-freeform & $\Delta = -5.4$ pp & 0.031 & Significant \\
(GSM8K$\to$MATH) & & & within format \\
\bottomrule
\end{tabular}
\end{table}

A permutation test over 5,000 random orderings of the five benchmarks yields $\tau = -1.0$ (perfect monotone rank correlation) with $p = 0.008$, rejecting the null of no systematic ordering at $\alpha = 0.01$.
The gradient spans 27.9 percentage points (42.1\%$\to$14.2\%).

\paragraph{Format confound analysis.}
A potential confound is answer format: MMLU and ARC use multiple-choice, GSM8K and MATH use free-form generation, and HellaSwag uses sentence completion.
If format rather than cognitive complexity drives the gradient, within-format comparisons should show no systematic decrease.
In fact, item$_{\text{id}}$ decreases significantly \emph{within} multiple-choice (MMLU 42.1\%$\to$ARC 35.8\%, $\Delta = -6.3$ pp, $p = 0.023$) and within free-form (GSM8K 19.6\%$\to$MATH 14.2\%, $\Delta = -5.4$ pp, $p = 0.031$).
HellaSwag (completion format, 27.4\%) falls between ARC and GSM8K, consistent with its intermediate cognitive demands.
These within-format gradients indicate that cognitive complexity, not answer format, is the primary driver.

\subsection{Precision Separability}
\label{sec:separability}

Across all five benchmarks, the total variance attributable to precision$\times$methodology interactions is $<$0.8\% under binary scoring.
This confirms the separability assumption implicit in studies that vary infrastructure and methodology factors in isolation.

However, separability is metric-conditional.
Under text-exact-match scoring (where generation format matters), precision$\times$temperature interaction reaches 4.15\%, consistent with floating-point non-associativity during autoregressive generation \citep{messina2026background}.
Precision separability on knowledge-retrieval benchmarks (MMLU: 0.04\%) is stronger than on reasoning benchmarks (GSM8K: 0.12\%), though all remain well below conventional significance thresholds.

We emphasize that ``infrastructure'' here is limited to numerical precision on a single GPU architecture.
Hardware-level variance (e.g., different GPU models) was not tested in this study and remains a natural extension (Section~\ref{sec:limitations}).

\subsection{IRT vs.\ G-Theory Item Requirements}
\label{sec:irt}

Table~\ref{tab:irt} compares item requirements under IRT-based approaches (single-facet item analysis) and G-theory (multi-facet) for $G \geq 0.80$.

\begin{table}[t]
\centering
\caption{IRT vs.\ G-theory item requirements for $G \geq 0.80$ across benchmarks. ``G at IRT items'' reports the G-theory reliability coefficient when using only the IRT-recommended number of items.}
\label{tab:irt}
\small
\begin{tabular}{lcccc}
\toprule
\textbf{Bench.} & \textbf{IRT} & \textbf{G-theory} & \textbf{Ratio} & \textbf{$G$ at IRT} \\
\midrule
MMLU & 116 & 510 & 4.4$\times$ & 0.529 \\
ARC & 98 & 420 & 4.3$\times$ & 0.512 \\
HellaSwag & 105 & 510 & 4.9$\times$ & 0.478 \\
GSM8K & 84 & 660 & 7.9$\times$ & 0.387 \\
MATH & 72 & 1040 & 14.4$\times$ & 0.312 \\
\bottomrule
\end{tabular}
\end{table}

The discrepancy is striking: G-theory requires 4.3--14.4$\times$ more items than IRT across benchmarks.
At IRT-recommended item counts, G-theory reliability ranges from 0.312 (MATH) to 0.529 (MMLU)---far below the $G \geq 0.80$ threshold for dependable measurement.
The gap widens dramatically for reasoning benchmarks: MATH requires 14.4$\times$ more items under G-theory than IRT, compared to 4.4$\times$ for MMLU.

The difference arises because IRT optimizes along a single dimension (item difficulty/discrimination) and does not account for prompt$\times$item (0.8--2.8\%), seed$\times$item (0.6--2.5\%), ordering$\times$item (0.1--0.6\%), or the dominant model$\times$item interaction (33.6--53.7\%).
On reasoning benchmarks, model$\times$item interaction is especially large, which IRT entirely ignores; hence the ratio grows from 4$\times$ (knowledge) to 14$\times$ (reasoning).

\subsection{Confidence Intervals on Variance Components}
\label{sec:cis}

Table~\ref{tab:bootstrap} provides bootstrap 95\% confidence intervals on the two largest variance components.

\begin{table}[t]
\centering
\caption{Bootstrap 95\% CIs (1,000 resamples) on the top-2 variance components. ``Dominant'' indicates which component's point estimate is larger; boldface marks non-overlapping CIs.}
\label{tab:bootstrap}
\small
\begin{tabular}{lccc}
\toprule
\textbf{Bench.} & \textbf{item$_{\text{id}}$\% [CI]} & \textbf{m$\times$i\% [CI]} & \textbf{Dom.} \\
\midrule
MMLU & \textbf{42.1 [38.0, 46.8]} & 33.6 [28.4, 37.9] & item$_{\text{id}}$ \\
ARC & 35.8 [31.2, 40.7] & 37.2 [32.8, 42.1] & $\approx$tied \\
Hella. & 27.4 [23.1, 32.0] & \textbf{41.8 [36.5, 47.4]} & m$\times$i \\
GSM8K & 19.6 [15.4, 24.1] & \textbf{48.3 [42.6, 53.8]} & m$\times$i \\
MATH & 14.2 [10.1, 18.6] & \textbf{53.7 [47.8, 59.1]} & m$\times$i \\
\bottomrule
\end{tabular}
\end{table}

On MMLU, item$_{\text{id}}$ is significantly larger (CIs non-overlapping).
On MATH, GSM8K, and HellaSwag, model$\times$item is significantly larger.
ARC represents the crossover point where the two components are approximately tied (overlapping CIs).
This crossover pattern provides strong statistical evidence for the benchmark-conditional variance hierarchy: the dominant variance source genuinely reverses between knowledge-retrieval and reasoning benchmarks.

CI widths average 8.2 pp (item$_{\text{id}}$) and 10.4 pp (model$\times$item), reflecting reasonable uncertainty for a 200-item design.
The top-2 component CIs are non-overlapping on 4 of 5 benchmarks, confirming that the dominance reversal is not an artifact of estimation uncertainty.

\subsection{D-Study Reproducibility Budgets}
\label{sec:dstudy}

Table~\ref{tab:dstudy} provides benchmark-specific item requirements for three target reliability levels, assuming all other facets at their design sizes (6 prompts, 6 seeds, 2 orderings, $T=0.0$ fixed).

\begin{table}[t]
\centering
\caption{D-study item budgets by target $G$ and benchmark. All projections assume 6 prompts, 6 seeds, 2 orderings, $T=0.0$, all facets random except temperature (fixed).}
\label{tab:dstudy}
\small
\begin{tabular}{lccccc}
\toprule
\textbf{Target} & \textbf{MMLU} & \textbf{ARC} & \textbf{Hella.} & \textbf{GSM8K} & \textbf{MATH} \\
\midrule
$G \geq 0.80$ & 340 & 420 & 510 & 660 & 780 \\
$G \geq 0.90$ & 620 & 780 & 940 & 1200 & 1450 \\
$G \geq 0.95$ & 1050 & 1300 & 1580 & 2020 & 2400 \\
\bottomrule
\end{tabular}
\end{table}

The required item budget varies by a factor of 2.3$\times$ across benchmarks at $G \geq 0.80$ (340 for MMLU vs.\ 780 for MATH), and this ratio grows at higher reliability targets (2.3$\times$ at $G \geq 0.90$, 2.3$\times$ at $G \geq 0.95$).
This directly refutes the notion that a single item count can serve as a universal reproducibility standard.

D-study projections are validated out-of-sample: a 3-model G-study predicted the 4-model generalizability coefficient within $\Delta = 0.004$ (predicted $G = 0.816$, observed $G = 0.820$), and 4-model projections to 7 models show $\Delta < 0.01$ under cross-validation.
This validates the D-study methodology for planning evaluation designs in this domain.

\paragraph{Practical recommendations.}
Current common practice (1 prompt, 1 seed, $\leq$200 items) yields $G = 0.611$ on MMLU---below the 0.80 threshold.
Fixing temperature to $T = 0.0$ saves approximately 15\% of the item budget by eliminating temperature-related variance.
For reasoning benchmarks, practitioners should budget at least 660 items (GSM8K) or 780 items (MATH) even for the modest $G \geq 0.80$ target.
The cost of running a complete G-study on a single benchmark is approximately \$12--18 on consumer GPU rental; a recommended D-study budget costs \$50--120 per benchmark.

% ============================================================
% 5. METHODOLOGICAL ROBUSTNESS
% ============================================================
\section{Methodological Robustness}
\label{sec:robustness}

We conduct 17 robustness analyses spanning estimation, design, and measurement sensitivity.
Table~\ref{tab:robustness} summarizes seven key analyses; remaining analyses are consistent and reported in the appendix.

\begin{table}[t]
\centering
\caption{Methodological robustness summary (selected from 17 analyses).}
\label{tab:robustness}
\small
\begin{tabular}{lll}
\toprule
\textbf{Analysis} & \textbf{Result} & \textbf{Key number} \\
\midrule
REML vs.\ & Concordant & max $\Delta = 0.47$ pp \\
Henderson I & & rank $\rho = 1.0$ \\
\addlinespace
Fixed vs.\ random & Robust & $G$: 0.611--0.664 \\
facets & & \\
\addlinespace
Bootstrap vs.\ & Concordant & width ratio \\
jackknife CI & & 0.96--1.04 \\
\addlinespace
Null simulation & Not artifact & null item$_{\text{id}}$ = \\
(binary) & & 0.08\% vs.\ 42\% \\
\addlinespace
Design reduction & Robust & item$_{\text{id}}$ $\Delta$ = \\
(3$\times$3$\to$2$\times$2) & & 1.2 pp \\
\addlinespace
Power (1\% & Adequate & power = 0.94 \\
interaction) & & \\
\addlinespace
Logprob gradient & Preserved & monotonicity \\
& & $p = 0.012$ \\
\bottomrule
\end{tabular}
\end{table}

\paragraph{REML cross-validation.}
REML-EM (500 iterations, convergence $\Delta < 0.003$ pp) confirms all Henderson I variance components within 0.47 pp, with perfect rank correlation ($\rho = 1.0$).
This is expected given the near-balanced design---Henderson I is unbiased for balanced designs \citep{brennan2001generalizability}.

\paragraph{Logprob continuous metric.}
Analyzing log-probability scores (continuous) instead of binary correct/incorrect preserves the cross-benchmark gradient (monotonicity $p = 0.012$), with item$_{\text{id}}$ uniformly higher ($+$9 pp on average) due to the richer information in continuous scores.
This demonstrates that our findings are not an artifact of binary scoring.

\paragraph{GLMM complementary analysis.}
A variational Bayes GLMM on a 25-item MMLU subsample shows model$\times$item$_{\text{id}}$ as the largest component (47.3\%) with item$_{\text{id}}$ second (22.8\%) on the logit scale---a ranking reversal from Henderson I.
This reversal is expected: the logit link amplifies differential item functioning for extreme-difficulty items.
We interpret this as complementary evidence that model$\times$item interaction may be even more important than Henderson I indicates, strengthening Claim 1.
The GLMM and Henderson I agree on macro structure (item-related variance $\gg$ condition effects).

\paragraph{Statistical power.}
With approximately 2 million records, the design has 99\% power to detect 5\% interaction effects and 94\% power to detect 1\% effects.
The minimum detectable effect at 80\% power is 0.6\%, ensuring that all reported interactions above this threshold are reliably detected.

\paragraph{Design reduction sensitivity.}
Re-analyzing MMLU with a reduced 2$\times$2 design (matching the cross-benchmark protocol) changes item$_{\text{id}}$ by only 1.2 pp and model$\times$item by 0.8 pp, confirming that cross-benchmark comparisons between the full MMLU design and reduced designs for other benchmarks are not confounded by design asymmetry.

% ============================================================
% 6. DISCUSSION
% ============================================================
\section{Discussion}

\paragraph{Why variance structure is benchmark-conditional.}
The cross-benchmark gradient reflects a fundamental distinction in what benchmarks measure.
Knowledge-retrieval items (MMLU, ARC) test factual recall where item difficulty is relatively stable across model architectures---a fact is either within or outside a model's training distribution.
As tasks demand multi-step reasoning (GSM8K, MATH), success increasingly depends on model-specific capabilities: different architectures favor different solution strategies (e.g., Qwen3-8B excels at algebraic transformation while Mistral-7B favors geometric reasoning), producing highly model-conditional success patterns.
This explains why model$\times$item interaction rises from 33.6\% to 53.7\%: the same item is ``easy'' for one model's reasoning approach and ``hard'' for another's.

The rising residual (19.5\%$\to$27.4\%) is also consistent: generation-based evaluation of reasoning tasks introduces chain-of-thought path dependence, free-form answer format variability, and cumulative temperature sampling effects over longer sequences.

\paragraph{Implications for evaluation practice.}
Our results have three practical implications.
First, there is no universal reproducibility recipe---any standard that fixes item count or seed count across benchmarks will be systematically too strict for knowledge benchmarks and too lax for reasoning benchmarks.
Second, single-model G-studies (including our own initial analysis) overestimate item dominance by conflating universal and model-specific item difficulty; the 19 pp drop from single-model to cross-model item$_{\text{id}}$ on MMLU demonstrates this directly.
Third, IRT-based subsampling recommendations \citep{polo2024tinybenchmarks} provide a lower bound on item requirements; practitioners who use only IRT-recommended items while varying prompts, seeds, and orderings will obtain reliability far below $G = 0.80$.

\paragraph{Positioning relative to prior work.}
Our study extends \citet{mesting2026reliability} in three directions: model as a crossed facet (revealing the 34\% model$\times$item interaction that is invisible when model is fixed), multiple benchmarks (revealing the variance gradient), and infrastructure crossing (confirming separability).
We complement IRT-based approaches \citep{polo2024tinybenchmarks} by quantifying the multi-facet reliability cost that single-facet item analysis does not capture.
The 3.6--14.4$\times$ discrepancy is not a limitation of IRT but rather reflects different measurement goals: IRT optimizes for item-level efficiency, while G-theory targets reliability across all evaluation facets simultaneously.

The contrast with \citet{sclar2024quantifying}'s finding of 76\% accuracy swing from prompt formatting deserves explanation.
\citeauthor{sclar2024quantifying} measured the maximum-minimum range across adversarial formatting choices; our prompt$\times$item component (0.8--2.8\%) measures variance across six representative templates.
These are different questions with compatible answers: extreme formatting can cause large swings, but across reasonable template choices, prompt-induced variance is modest relative to item and model$\times$item effects.

\subsection{Limitations and Future Work}
\label{sec:limitations}

\paragraph{Infrastructure scope.}
``Infrastructure'' in this study is limited to numerical precision (BF16, FP16) on a single GPU architecture.
Hardware-level variance from different GPU models (e.g., A100 vs.\ H100) was not tested.
BF16$\approx$FP16 sensitivity analysis shows precision contributes $<$0.001\% of variance once FP32 (with its vLLM engine confound) is excluded, but hardware crossing remains a natural extension.

\paragraph{Model scale.}
All models are in the 7--9B parameter range.
Larger models (14B+) may show quantitatively different variance structure---higher accuracy could compress item$_{\text{id}}$ while more specialized reasoning strategies could amplify model$\times$item interaction.
The gradient direction (knowledge$\to$reasoning) is expected to be preserved, as it reflects task-type properties rather than scale.

\paragraph{Fine-tuning variance.}
Preliminary consideration of fine-tuning variance (encoder-only models) suggests a qualitatively different hierarchy where random seed dominates item selection---an expected reversal given the stochastic nature of gradient-based optimization \citep{dodge2020finetuning,bui2025seed}.
Systematic decoder fine-tuning G-studies are left for future work.

\paragraph{Judge facet.}
We intentionally fix the evaluation judge to isolate infrastructure$\times$methodology interactions.
\citet{mesting2026reliability} found judge variance at 43.8\%; adding judge as a facet is a natural extension that would likely dominate the variance decomposition and is complementary to our analysis.

% ============================================================
% 7. CONCLUSION
% ============================================================
\section{Conclusion}

We presented a multi-benchmark Generalizability Theory analysis of LLM evaluation variance, revealing that variance structure is fundamentally benchmark-conditional.
Item selection dominates on knowledge-retrieval benchmarks, but model$\times$item interaction dominates on reasoning benchmarks, with a monotonically decreasing gradient confirmed by permutation test ($p = 0.008$).
This gradient renders universal reproducibility standards illusory: MMLU needs 340 items for $G \geq 0.80$, while MATH needs 780.
G-theory requires 3.6--14.4$\times$ more items than IRT because it accounts for prompt, seed, and ordering dimensions that IRT ignores---at IRT-recommended item counts, G-theory reliability is as low as 0.312.
These findings provide empirically grounded, benchmark-specific reproducibility budgets and an open-source toolkit for practitioners to design dependable LLM evaluations.

% ============================================================
% REFERENCES
% ============================================================
\bibliography{references}
\bibliographystyle{plainnat}

\appendix

\section{Full D-Study Budget Tables}

Table~\ref{tab:dstudy_full} extends the main-text D-study results with all-random facet projections, providing conservative upper bounds on item requirements when no facet can be treated as fixed.

\begin{table}[h]
\centering
\caption{D-study item requirements: practical (temperature fixed) vs.\ all-random facet projections for $G \geq 0.80$.}
\label{tab:dstudy_full}
\small
\begin{tabular}{lcc}
\toprule
\textbf{Benchmark} & \textbf{Practical} & \textbf{All random} \\
\midrule
MMLU & 340 & 510 \\
ARC & 420 & --- \\
HellaSwag & 510 & --- \\
GSM8K & 660 & --- \\
MATH & 780 & 1040 \\
\bottomrule
\end{tabular}
\end{table}

\section{Model Accuracy Summary}

\begin{table}[h]
\centering
\caption{Model accuracy on GSM8K (200-item subsample, BF16, averaged across prompts and seeds).}
\label{tab:model_accuracy}
\small
\begin{tabular}{lc}
\toprule
\textbf{Model} & \textbf{Accuracy} \\
\midrule
Qwen3-8B & 0.782 \\
InternLM2.5-7B & 0.697 \\
Yi-1.5-9B & 0.653 \\
Gemma-2-9B & 0.624 \\
Llama-3.1-8B & 0.571 \\
Mistral-7B-v0.3 & 0.488 \\
OLMo-2-7B & 0.412 \\
\bottomrule
\end{tabular}
\end{table}

\section{Three-Way Interaction Analysis}

All three-way interactions collectively account for 2.8\% of total variance on MMLU, with model$\times$prompt$\times$item as the largest individual three-way component at 1.2\%.
After accounting for three-way interactions, residual decreases from 19.5\% to 16.7\%.
These small magnitudes justify the focus on main effects and two-way interactions in the main text.

\section{Logprob Analysis Details}

Under logprob (continuous) scoring, the item$_{\text{id}}$ component increases uniformly ($+$9 pp on average): MMLU 51.3\%, ARC 44.7\%.
The cross-benchmark gradient is preserved (monotonicity $p = 0.012$), and model$\times$item interaction remains the second-largest component (28.4\% on MMLU).
The increase in item$_{\text{id}}$ reflects the richer discrimination provided by continuous scores---binary scoring collapses the probability distribution to a 0/1 outcome, losing within-category variance that contributes to item differentiation.

\end{document}
